\subsection{Sprint 0}

Durante a sprint inicial, o foco principal recaiu sobre a integração da equipe e a definição de uma base sólida de design e arquitetura, atividades que visavam estabelecer a base para o desenvolvimento e estruturação do projeto SalesPilot. Ao longo desse período, busquei liderar ativamente a elaboração do \textit{style guide} e a prototipação no Figma, assumindo a responsabilidade de ensinar os colegas de time a utilizar a ferramenta e a compreender a importância da componentização para a escalabilidade do projeto. Já buscando exercer meu papel de \ac{ages} III, procurei adiantar a escolha da biblioteca de componentes Material UI (MUI) \cite{MUI}, para que pudéssemos otimizar o processo produtivo ao aproveitar elementos da \textit{library} criados e disponibilizados gratuitamente pela comunidade, garantindo que o protótipo fosse fácil de implementar e reproduzir tecnicamente depois.

Durante os encontros com os stakeholders, o time promoveu momentos de questionamento para tentar sanar a falta de clareza que ainda permeava o escopo do projeto. Percebeu-se que a estrutura do banco de dados e as informações específicas que os clientes desejavam visualizar na plataforma, principalmente no painel geral de cada \textit{role}, ainda pareciam incertos e pouco definidos. Essa indefinição exigiu diversas discussões internas com a equipe, nas quais buscamos alinhar nosso entendimento sobre o projeto como um todo e mitigar os riscos de desenvolvermos funcionalidades que não atendessem plenamente às necessidades de administradores, gestores e vendedores ou que não cumprissem com as expectativas iniciais dos \textit{stakeholders}.

Apesar de ter conseguido avançar no suporte aos alunos de \ac{ages} I e adiantar a estruturação inicial do repositório de frontend com React 19 \cite{React}, Vite \cite{Vite} e MUI \cite{MUI} já definidos, enfrentamos obstáculos significativos como time. A ausência de um calendário rígido e a indefinição de tasks claras no início da disciplina, somadas a imprevistos de horários e eventos, dificultaram uma progressão mais linear. Assim, essa experiência reforçou a lição de que o planejamento deve ser ágil para mitigar mudanças de escopo, requisitos vagos ou imprecisos e conflitos inesperados de agenda. Como próximos passos, focaremos em estreitar as definições de dados com os clientes para solidificar a arquitetura e iniciar formalmente as tarefas de implementação dos repositórios e ambientes de desenvolvimento.